\section{System Architecture}
\label{sec:architecture}

CircuitCraft is built on Fabric, a lightweight, widely used modding toolchain for Minecraft.
The codebase is organized into five packages, shown in Figure~\ref{fig:architecture}: a
solver package with no dependency on any Minecraft class at all, a pair of packages
describing the in-world blocks and their per-block simulation state, a networking/topology
package that ties placed blocks together into a circuit and synchronizes probe readings to
clients, and a client-only package rendering the oscilloscope heads-up display.

\begin{figure}[htbp]
\centering
\begin{tikzpicture}[
    box/.style={draw, rounded corners, minimum width=6.2cm, minimum height=1cm,
                align=center, font=\small},
    node distance=0.4cm
]
\node[box, fill=gray!10] (sim) {\texttt{sim} --- pure-Java MNA circuit solver\\(no Minecraft dependency)};
\node[box, below=of sim, fill=blue!8] (blockentity) {\texttt{blockentity} --- per-block simulation state\\(components, wire/ground nodes, modules)};
\node[box, below=of blockentity, fill=blue!8] (block) {\texttt{block} --- placement, orientation, right-click interaction};
\node[box, below=of block, fill=orange!12] (network) {\texttt{network} --- topology construction, solve loop,\\multi-channel probe protocol};
\node[box, below=of network, fill=green!10] (client) {\texttt{client} --- oscilloscope HUD rendering (client-side only)};

\draw[-{Latex[length=2mm]}] (blockentity) -- (sim);
\draw[-{Latex[length=2mm]}] (block) -- (blockentity);
\draw[-{Latex[length=2mm]}] (network) -- (blockentity);
% Routed as a dogleg out to the right, rather than a straight line, since network and sim
% are two boxes apart in the stack -- a direct line between their centers/anchors would cut
% straight through the block and blockentity boxes' text in between.
\draw[-{Latex[length=2mm]}] (network.east) -- ++(0.6,0) |- (sim.east);
\draw[-{Latex[length=2mm]}] (client) -- (network);
\end{tikzpicture}
\caption{Package-level architecture. The solver has zero Minecraft dependencies and is
directly unit-testable with plain \texttt{javac}/\texttt{java}.}
\label{fig:architecture}
\end{figure}

A two-terminal component occupies a single block, oriented along the horizontal or vertical
axis the player was facing when placed; its two electrical leads are its front and back
faces along that axis, and its other four faces are electrically insulated. Once per game
tick, for each independent group of adjacent, mutually-conductive blocks (a ``network''), the
mod rebuilds a union--find partition over every block face presenting a conductive surface
toward a conductive neighbor, assigns one MNA node index to each resulting partition, and
instructs every component in that network to stamp itself into the circuit between its
resolved node indices. The mod provides seventeen such blocks and items in total, grouped by
electrical role: six circuit elements (resistor, capacitor, inductor, memristor, diode, and
an ideal op-amp -- the mod's only three-terminal component, its input-pair keyed as a
separate lead from its output without any per-terminal-count branching in the algorithm
above); two purely topological blocks (a wire, conductive on all six faces with no
orientation, and a ground block, identical to wire for topology but additionally
establishing the network's voltage reference); three independent voltage sources (a
fixed-DC power supply, a function generator producing sine/square/triangle waves whose
amplitude and frequency are overridden by two attachable modules, and a dedicated AC Source
used only by the Bode-plot probe below, Section~\ref{sec:ac-analysis}) -- all three require a
redstone signal before driving anything (Section~\ref{sec:pedagogy}); and four measurement
tools (an in-circuit ammeter, electrically an ideal 0~V source reporting an exact branch
current, plus three handheld probes described next).

A handheld probe item lets a player pin up to three components, wires, or ground blocks
simultaneously; the server streams each pinned position's instantaneous voltage, current, a
rolling history, and a summary once per tick, rendered as independently auto-scaled
oscilloscope traces (Figure~\ref{fig:scope}). A second probe item pins exactly two channels
and plots one against the other (an X-Y, Lissajous-style display) rather than against time,
each axis independently scaled to its own channel's peak magnitude. A third, AC probe
(Section~\ref{sec:ac-analysis}) is a genuinely two-step interaction: pinning an AC Source
first designates it as the sweep's excitation, and pinning a second point afterward runs the
sweep and displays the resulting Bode plot.
